% GENERATED FILE. Edit canonical lessons, then run npm run generate:books.
% canonical-source-hash: fnv1a64:41df87e43aa6e964
% canonical-lessons: RU-C379-zapasnoy, RU-C379-v-lyubom-sluchae, RU-C379-kolebatsya, RU-C379-reshatsya, RU-C379-v-zavisimosti-ot

\chapter{A fallback plan}
\label{ch:ru-a-fallback-plan}
\hlchaptermodality{379}

\begin{chapteropening}
\textbf{By the end of this chapter:} \emph{I can have a fallback ready, say in any case and depending on, hesitate, and make up my mind.}

Complete the last lesson of chapter 379: I can have a fallback ready, say in any case and depending on, hesitate, and make up my mind.
\end{chapteropening}

\section[\texorpdfstring{\ru{запасной} (zapasnóy)}{запасной (zapasnóy)}]{\ru{запасной} (zapasnóy) — spare, reserve}

\label{lesson:RU-C379-zapasnoy}



\begin{quote}
Before the new one: say the Russian for an alternative, then the Russian for final.
\end{quote}

\subsection*{You'll want to know: \ru{запасной}}
\textbf{\ru{запасной}} (\emph{zapasnóy}) — \textquotedblleft{}spare, reserve\textquotedblright{}.

\begin{quote}
\ru{У} \ru{меня} \ru{есть} \ru{запасной} \ru{ключ}.
\end{quote}

\begin{quote}
\emph{U menyá yest zapasnóy klyuch.}
\end{quote}

\begin{quote}
I have a spare key.
\end{quote}

\emph{\ru{Запасной} \ru{выход}} — an emergency exit. \emph{\ru{Запасной} \ru{план}} — a fallback plan.

\begin{cousinweb}
\textbf{\ru{запасной}} is built on \emph{\ru{запас}} \textquotedblleft{}a store, a supply\textquotedblright{}, which is usually traced to \emph{\ru{пасти}} \textquotedblleft{}to guard, to keep\textquotedblright{}: what is kept in reserve.
\end{cousinweb}

\begin{grammarlens}[title={What you've built}]
One more word for telling, explaining and describing at B1.
\end{grammarlens}

\subsection*{Your turn}
\begin{itemize}
\raggedright
  \item \emph{Say it:} \emph{zapasnóy}
  \item \emph{Say it:} \emph{zapasnóy}, once more
  \item \emph{Say it:} a sentence of your own with \emph{zapasnóy}
  \item \emph{Recall:} say the Russian for an alternative, then the Russian for final, then say \emph{zapasnóy} again
\end{itemize}

\begin{tcolorbox}[breakable,colback=teal!4,colframe=teal!35!black,title={Before you move on}]
What does \ru{запасной} mean? (\textquotedblleft{}Spare, reserve\textquotedblright{}.) Say it once more.
\end{tcolorbox}
\section[\texorpdfstring{\ru{в} \ru{любом} \ru{случае} (v lyubóm slúchaye)}{в любом случае (v lyubóm slúchaye)}]{\ru{в} \ru{любом} \ru{случае} (v lyubóm slúchaye) — in any case, anyway}

\label{lesson:RU-C379-v-lyubom-sluchae}



\begin{quote}
Before the new one: say the Russian for final, then the Russian for spare.
\end{quote}

\subsection*{You'll want to know: \ru{в} \ru{любом} \ru{случае}}
\textbf{\ru{в} \ru{любом} \ru{случае}} (\emph{v lyubóm slúchaye}) — \textquotedblleft{}in any case, anyway\textquotedblright{}.

\begin{quote}
\ru{В} \ru{любом} \ru{случае} \ru{я} \ru{тебе} \ru{позвоню}.
\end{quote}

\begin{quote}
\emph{V lyubóm slúchaye ya tebé pozvonyú.}
\end{quote}

\begin{quote}
In any case I will call you.
\end{quote}

\emph{\ru{Любой}} means any. It closes an argument: whatever happens.

\begin{cousinweb}
\textbf{\ru{в} \ru{любом} \ru{случае}} is \textquotedblleft{}in any case\textquotedblright{}: \emph{\ru{любой}} \textquotedblleft{}any\textquotedblright{}, on the root of \emph{\ru{любить}} \textquotedblleft{}to love\textquotedblright{}, and \emph{\ru{случай}} \textquotedblleft{}a case\textquotedblright{}, from Proto-Slavic \emph{*sŭlučajĭ} (EtymDB).
\end{cousinweb}

\begin{grammarlens}[title={What you've built}]
One more word for telling, explaining and describing at B1.
\end{grammarlens}

\subsection*{Your turn}
\begin{itemize}
\raggedright
  \item \emph{Say it:} \emph{v lyubóm slúchaye}
  \item \emph{Say it:} \emph{v lyubóm slúchaye}, once more
  \item \emph{Say it:} a sentence of your own with \emph{v lyubóm slúchaye}
  \item \emph{Recall:} say the Russian for final, then the Russian for spare, then say \emph{v lyubóm slúchaye} again
\end{itemize}

\begin{tcolorbox}[breakable,colback=teal!4,colframe=teal!35!black,title={Before you move on}]
What does \ru{в} \ru{любом} \ru{случае} mean? (\textquotedblleft{}In any case, anyway\textquotedblright{}.) Say it once more.
\end{tcolorbox}
\section[\texorpdfstring{\ru{колебаться} (kolebátsya)}{колебаться (kolebátsya)}]{\ru{колебаться} (kolebátsya) — to hesitate, to waver}

\label{lesson:RU-C379-kolebatsya}



\begin{quote}
Before the new one: say the Russian for spare, then the Russian for in any case.
\end{quote}

\subsection*{You'll want to know: \ru{колебаться}}
\textbf{\ru{колебаться}} (\emph{kolebátsya}) — \textquotedblleft{}to hesitate, to waver\textquotedblright{}.

\begin{quote}
\ru{Я} \ru{долго} \ru{колебался}, \ru{но} \ru{потом} \ru{согласился}.
\end{quote}

\begin{quote}
\emph{Ya dólgo kolebálsya, no potóm soglasílsya.}
\end{quote}

\begin{quote}
I hesitated for a long time, but then I agreed.
\end{quote}

Imperfective; present \emph{\ru{колеблюсь}, \ru{колеблешься}}. Literally to sway: \emph{\ru{пламя} \ru{колеблется}}.

\begin{cousinweb}
\textbf{\ru{колебаться}} is inherited from Proto-Slavic \emph{*kolěbati} (EtymDB), usually glossed \textquotedblleft{}to rock, to sway\textquotedblright{}; a person who hesitates sways from one side to the other.
\end{cousinweb}

\begin{grammarlens}[title={What you've built}]
One more word for telling, explaining and describing at B1.
\end{grammarlens}

\subsection*{Your turn}
\begin{itemize}
\raggedright
  \item \emph{Say it:} \emph{kolebátsya}
  \item \emph{Say it:} \emph{kolebátsya}, once more
  \item \emph{Say it:} a sentence of your own with \emph{kolebátsya}
  \item \emph{Recall:} say the Russian for spare, then the Russian for in any case, then say \emph{kolebátsya} again
\end{itemize}

\begin{tcolorbox}[breakable,colback=teal!4,colframe=teal!35!black,title={Before you move on}]
What does \ru{колебаться} mean? (\textquotedblleft{}To hesitate, to waver\textquotedblright{}.) Say it once more.
\end{tcolorbox}
\section[\texorpdfstring{\ru{решаться} (reshátsya)}{решаться (reshátsya)}]{\ru{решаться} (reshátsya) — to make up one’s mind, to dare}

\label{lesson:RU-C379-reshatsya}



\begin{quote}
Before the new one: say the Russian for in any case, then the Russian for to hesitate.
\end{quote}

\subsection*{You'll want to know: \ru{решаться}}
\textbf{\ru{решаться}} (\emph{reshátsya}) — \textquotedblleft{}to make up one’s mind, to dare\textquotedblright{}.

\begin{quote}
\ru{Он} \ru{никак} \ru{не} \ru{решается} \ru{сказать} \ru{ей} \ru{правду}.
\end{quote}

\begin{quote}
\emph{On nikák ne resháyetsya skazát yey právdu.}
\end{quote}

\begin{quote}
He cannot bring himself to tell her the truth.
\end{quote}

Imperfective; perfective \emph{\ru{решиться}}. \emph{\ru{Решаться} \ru{на}} + accusative, or + infinitive. Not the same as \emph{\ru{решать}}, to solve or decide.

\begin{cousinweb}
\textbf{\ru{решаться}} is the reflexive of \emph{\ru{решать}} \textquotedblleft{}to decide\textquotedblright{}, inherited from Proto-Slavic \emph{*rěšiti} \textquotedblleft{}to untie, to loosen\textquotedblright{}: to make up your mind is to untie yourself from doubt.
\end{cousinweb}

\begin{grammarlens}[title={What you've built}]
One more word for telling, explaining and describing at B1.
\end{grammarlens}

\subsection*{Your turn}
\begin{itemize}
\raggedright
  \item \emph{Say it:} \emph{reshátsya}
  \item \emph{Say it:} \emph{reshátsya}, once more
  \item \emph{Say it:} a sentence of your own with \emph{reshátsya}
  \item \emph{Recall:} say the Russian for in any case, then the Russian for to hesitate, then say \emph{reshátsya} again
\end{itemize}

\begin{tcolorbox}[breakable,colback=teal!4,colframe=teal!35!black,title={Before you move on}]
What does \ru{решаться} mean? (\textquotedblleft{}To make up one’s mind, to dare\textquotedblright{}.) Say it once more.
\end{tcolorbox}
\section[\texorpdfstring{\ru{в} \ru{зависимости} \ru{от} (v zavísimosti ot)}{в зависимости от (v zavísimosti ot)}]{\ru{в} \ru{зависимости} \ru{от} (v zavísimosti ot) — depending on}

\label{lesson:RU-C379-v-zavisimosti-ot}



\begin{quote}
Before the new one: say the Russian for to hesitate, then the Russian for to make up one’s mind.
\end{quote}

\subsection*{You'll want to know: \ru{в} \ru{зависимости} \ru{от}}
\textbf{\ru{в} \ru{зависимости} \ru{от}} (\emph{v zavísimosti ot}) — \textquotedblleft{}depending on\textquotedblright{}.

\begin{quote}
\ru{В} \ru{зависимости} \ru{от} \ru{погоды} \ru{мы} \ru{поедем} \ru{на} \ru{море} \ru{или} \ru{в} \ru{горы}.
\end{quote}

\begin{quote}
\emph{V zavísimosti ot pogódy my poyédem na móre íli v góry.}
\end{quote}

\begin{quote}
Depending on the weather, we will go to the sea or to the mountains.
\end{quote}

\emph{\ru{От}} + genitive: \emph{\ru{в} \ru{зависимости} \ru{от} \ru{цены}}.

\begin{cousinweb}
\textbf{\ru{в} \ru{зависимости} \ru{от}} is built on \emph{\ru{зависимость}} \textquotedblleft{}dependence\textquotedblright{}, from \emph{\ru{зависеть}} \textquotedblleft{}to depend\textquotedblright{}, made of \emph{\ru{за}-} and \emph{\ru{висеть}} \textquotedblleft{}to hang\textquotedblright{}: one thing hangs on another.
\end{cousinweb}

\begin{grammarlens}[title={What you've built}]
One more word for telling, explaining and describing at B1.
\end{grammarlens}

\subsection*{Your turn}
\begin{itemize}
\raggedright
  \item \emph{Say it:} \emph{v zavísimosti ot}
  \item \emph{Say it:} \emph{v zavísimosti ot}, once more
  \item \emph{Say it:} a sentence of your own with \emph{v zavísimosti ot}
  \item \emph{Recall:} say the Russian for to hesitate, then the Russian for to make up one’s mind, then say \emph{v zavísimosti ot} again
\end{itemize}

\begin{tcolorbox}[breakable,colback=teal!4,colframe=teal!35!black,title={Before you move on}]
What does \ru{в} \ru{зависимости} \ru{от} mean? (\textquotedblleft{}Depending on\textquotedblright{}.) Say it once more.
\end{tcolorbox}
